%% maestro2tex -- generated table; do not edit by hand.
\begin{table}[htbp]
  \centering
  \caption[{Control-loop stability driving an electrochemical cell, by process corner.}]{Control-loop stability driving an electrochemical cell, by process corner. Cell-drive bench (Figure~\ref{fig:tb-celldrive}): Randles cell with $R_{\mathrm{CE}} = \SI{1}{\kilo\ohm}$, $R_u = \SI{1}{\kilo\ohm}$, $R_{\mathrm{ct}} = \SI{10}{\kilo\ohm}$, $C_{\mathrm{dl}} = \SI{1}{\micro\farad}$, working electrode at \SI{1.25}{\volt}, $C_L = \SI{5}{\pico\farad}$, $C_c = \SI{500}{\femto\farad}$, \SI{2.5}{\volt} supply, \SI{37}{\celsius} at the typical point and \SI{25}{\celsius} at the four skew corners. The probe sits in the reference-electrode sense wire, so the loop measured here includes the cell; it is a different loop from Table~\ref{tab:ampab-unity} and the two are not comparable. Loop gain is \SI{10}{\decibel} lower than the unity-feedback figure because the cell's \SI{12}{\kilo\ohm} loads an open-loop output resistance of about \SI{25}{\kilo\ohm}; crossover halves with it and the phase margin improves. These corner spreads are a lower bound: the corner set holds the passive model rows at typical (pinned-passive caveat, page~\pageref{p:pinned-passives}).}
  \label{tab:ampab-cellstab}
  \begin{tabular}{lrrrrrrr}
    \toprule
    Parameter & Unit & tt & ff & fs & sf & ss & Spread \\
    \midrule
    Loop gain at DC & \si{\decibel} & 72.24 & 71.85 & 72.24 & 73.57 & 74.11 & 2.26 \\
    Loop crossover & \si{\mega\hertz} & 10.440 & 10.835 & 10.508 & 10.533 & 10.223 & 0.612 \\
    Phase margin & \si{\degree} & 81.88 & 82.14 & 82.05 & 81.83 & 81.68 & 0.45 \\
    \bottomrule
  \end{tabular}
\end{table}
